% from ega4-20-fr.tex
\oldpage[IV-4]{248}
coïncident dans l'ouvert partout dense $U\cap W^{(x')}\cap W^{(x'')}$ de cet ensemble.
Il y a donc un morphisme $h_i:U\cup U_i\to Y$ qui prolonge $h$ à un ouvert $U\cup U_i$, contenant $E_i$ ; en outre, pour tout couple d'indices $i$, $j$, les graphes des restrictions $h_i|(U_i\cap U_j)$ et $h_j|(U_i\cap U_j)$ sont contenus dans $V_i\cap V_j$ ; comme $V_i\cap V_j$ est séparé sur $X$ et que les morphismes précédents coïncident dans un ouvert partout dense $U\cap U_i\cap U_j$ de $U_i\cap U_j$, ils sont égaux.
Le morphisme $h'$ égal à $h$ dans $U$, à $h_i$ dans chacun des $U_i$, répond à la question.
\end{proof}

\begin{remark}[20.4.8]
\label{IV.20.4.8}
Lorsque $E=X-U$, il est clair que si, pour tout ouvert affine $T$ de $X$, il existe un $S$-morphisme $h'_T:T\to Y$ prolongeant $h|(U\cap T)$ et tel que le composé $\Spec(k(x))\to T\xrightarrow{h'_T}Y$ soit égal à $h'_x$ pour tout $x\in T-(U\cap T)$, alors les $h'_T$ sont les restrictions d'un $S$-morphisme $h':X\to Y$ (partout défini) en vertu de l'assertion d'unicité dans \sref{IV.20.4.6}.
Pour prouver l'existence de $h'$, on est donc ramené au cas où $X$ est \emph{affine}, et alors il suffit que l'ensemble $F$, réunion des graphes des $h'_x$, soit \emph{quasi-compact} dans $X\times_S Y$ pour que l'on puisse appliquer \sref{IV.20.4.7}.
Il en sera ainsi lorsque les $h'_x$ sont de la forme $\Spec(k(x))\to Z\xrightarrow{h_0}Y$, où $Z$ est un sous-préschéma fermé de $X$ ayant $X-U$ pour espace sous-jacent, et $h_0$ un $S$-morphisme.
\end{remark}

\begin{corollary}[20.4.9]
\label{IV.20.4.9}
Soient $S$ un préschéma localement noethérien, $X$ un préschéma localement noethérien, $f:X\to S$ un morphisme plat, $g:Y\to S$ un morphisme localement de type fini.
Soient $U$ un ouvert dense dans $S$, $h:f^{-1}(U)\to Y$ un $S$-morphisme, $T=S-U$, $Z$ le sous-préschéma fermé réduit de $X$ ayant $f^{-1}(T)$ comme espace sous-jacent, $h_0:Z\to Y$ un $S$-morphisme.
Supposons $X$ normal en tous les points de $Z$.
Afin qu'il existe un $S$-morphisme (nécessairement unique) $h':X\to Y$ prolongeant $h$ et $h_0$, il faut et il suffit que la condition suivante soit satisfaite :

\emph{Pour tout spectre $S_1$ d'un anneau de valuation discrète, de point fermé $a$ et de point générique $b$, et tout morphisme $u:S_1\to S$ tel que $u(a)\in T$ et $u(b)\in U$, il existe un $S_1$-morphisme $h'_1:X_{(S_1)}\to Y_{(S_1)}$ qui prolonge $h_{(S_1)}:f_{(S_1)}^{-1}(b)\to Y_{(S_1)}$ et est tel que le diagramme}
\[
  \xymatrix{
  \Spec(k(z_1)) \ar[r] \ar[d] & Z_{(S_1)} \ar[d]^{(h_0)_{(S_1)}} \\
  X_{(S_1)} \ar[r]_{h'_1} & Y_{(S_1)}
  }
\]
\emph{soit commutatif pour tout $z_1\in Z_{(S_1)}$.}

L'hypothèse que $f$ est plat entraîne en effet que $f^{-1}(U)$ est dense dans $X$ \sref{IV.2.3.10}, et il suffit alors d'appliquer \sref{IV.20.4.8}.
\end{corollary}

\begin{corollary}[20.4.10]
\label{IV.20.4.10}
Sous les hypothèses de \sref{IV.20.4.6}, supposons en outre $Y$ séparé et localement quasi-fini sur $S$.
Soient $U$ un ouvert dense dans $X$, $h:U\to Y$ un $S$-morphisme, $\omega$ la $S$-application rationnelle correspondante, $x\in X-U$ un point normal de $X$.
Afin que $\omega$ soit définie au point $x$, il faut et il suffit que la condition suivante soit vérifiée : il existe un spectre $S_1$ d'un anneau de valuation discrète, de point fermé $a$ et de point générique $b$, et un morphisme $u:S_1\to X$
